\section{A center adapted to the multiplicity defects}
\label{sec:centering}

We now use the signed estimate to construct a single reference graph.  Its
height error will be supported in a set of vanishing density.  Two features
of the construction are needed later: neighboring minimal fits agree to
\emph{second} order in the height error, and failure of a fit forces a
definite amount of lower multiplicity in every fixed interior subball.

Minimal-graph refinement and Jacobi regularization of an average appear in
Brena--De Lellis--Franceschini \cite[Sections~2 and~6.3--6.4]{BDLF2025}.
Almgren developed the center-manifold approach to higher-multiplicity
regularity \cite{Almgren2000}. De Lellis--Spadaro construct a center
manifold from smoothed local averages and Whitney gluing
\cite[Sections~1 and~4]{DLS2016II}. Here the fits are selected using
multiplicity defects, and we prove their quadratic compatibility before
gluing them.

\subsection{Minimal fits and separation}

We specify the points at which the construction is made.  Translate such a
point to the origin, and write $P_0=\R^m\times\{0\}$ and $\pi_0$ for its
orthogonal projection.  Define the dilated varifold explicitly by
$V_r(\Phi)=r^{-m}\int\Phi(X/r,T_XE)\,\dd\mu(X)$.
The hypotheses are
\begin{equation}
 \label{ce:typical}
 V_r\longrightarrow Q|P_0|,
 \qquad
 r^{-m}\mu\bigl(\mathbf B_r(0)\cap\{\theta\ne Q\}\bigr)
       \longrightarrow0,
 \qquad Q\in\N\setminus\{0\},
\end{equation}
where the first convergence includes the tangent-plane variable.
Thus dilated masses and tangent projections converge to those of the
multiplicity-$Q$ plane.  As usual, $\theta$ is the integer multiplicity
at almost every point of the rectifiable set.

These hypotheses also give a support statement.  In a sufficiently small
punctured ambient ball,
\begin{equation}
 \label{ce:cone}
 |X_v|\le\varepsilon(|X|)|X_h|\quad(X\in\spt\mu),
 \qquad \varepsilon(s)\longrightarrow0.
\end{equation}
Indeed, if $|X_v|\ge\varepsilon_0|X|$ along a sequence $X\to0$, the balls
$\mathbf B_{c\varepsilon_0|X|}(X)$ have mass at least
$c'\varepsilon_0^m|X|^m$ by \Cref{lt:monotonicity}.  After dilation by
$|X|^{-1}$ these balls stay a fixed positive distance from $P_0$, contrary
to \eqref{ce:typical}.  Passing from $|X|$ to $|X_h|$ changes the function
$\varepsilon$ only by a factor tending to one.

Choose a cylinder strictly inside that ambient ball, with the support
separated from its vertical boundary, and dilate it to contain the fixed
base ball $B_{300}\subset P_0$.  All varifolds below are restrictions to
this cylinder, and all tests have horizontal support away from its side.
The mass of each fixed interior ball is bounded above by a constant
depending on the fixed data, including $Q$, times its radius to the power $m$: apply monotonicity at that
center and a fixed larger radius.  By choosing the initial dilation
sufficiently small, the height on the fixed cylinder and the weak mass
error can be made arbitrarily small.  In addition the measure
\begin{equation}
 \label{ce:defectmeasure}
 \lambda=(\pi_0)_\#\bigl(\mu\llcorner\{\theta\ne Q\}\bigr)
\end{equation}
satisfies, for all radii needed in the fixed chart,
\begin{equation}
 \label{ce:initialdefect}
 \lambda(B_s(0))\le\eta s^m,
 \qquad \lambda(B_s(0))=o(s^m)\quad(s\downarrow0).
\end{equation}
Here $\eta>0$ can be prescribed before making the dilation.  To see this,
\eqref{ce:cone} puts the relevant part of the cylinder above $B_s(0)$
inside $\mathbf B_{2s}(0)$, and then one uses the second limit in
\eqref{ce:typical}.  We will choose one sufficiently small $\eta$ below
and keep the resulting data fixed.

\begin{lemma}[Minimal fits and their compatibility]
 \label{ce:fit}
Fix a finite derivative order $J\ge6$.  There are an integer
$k_0>J+4$, integers $q,N$, and constants $C,m_*>0$, depending only on the
dimensions, $Q$, and these fixed orders, with the following property.
For $0<m_0<m_*$ and sufficiently small initial height and mass error,
one can make the following construction simultaneously for every
$x\in B_5$.

At the scales $\ell_j=2^{-j}$ construct smooth minimal graphs
$S_{x,j}=\operatorname{graph}f_{x,j}$ over $B_{80\ell_j}(x)$, and stop at
the first scale for which
\begin{equation}
 \label{ce:accept}
 \sup\bigl\{|X_v-f_{x,j}(X_h)|:
       X\in\spt\mu,\ X_h\in B_{64\ell_j}(x)\bigr\}
       \le m_0\ell_j^{k_0}
\end{equation}
fails.  If there is no failure, all scales are constructed.  The graph at
the first failed scale is included in the construction.  On fixed
interior enlargements of its test ball, every constructed graph satisfies
the height upper bound $C m_0\ell_j^{k_0}$, the tilt and mass-defect bounds
\begin{equation}
 \label{ce:fitenergy}
 \ell_j^{-m}\bigl(D+|M|\bigr)(B_{70\ell_j}(x))
       \le C a_j^2,
 \qquad a_j=m_0\ell_j^{k_0-1},
\end{equation}
and the following centering estimate.  For each smooth vector test
$\zeta$ supported in $B_{70}$,
\begin{equation}
 \label{ce:moment}
 \left|\ell_j^{-m}\int
   \frac{X_v-f_{x,j}(X_h)}{\ell_j}\cdot
   \zeta\left(\frac{X_h-x}{\ell_j}\right)\,\dd\mu(X)\right|
       \le C a_j^2\|\zeta\|_{H^q}.
\end{equation}
The same assertion holds in normal coordinates over $S_{x,j}$, with
base volume in the definition of the distribution.

Every graph has physical derivatives of orders $1,\ldots,J$ bounded by
$Cm_0$ in the fixed initial coordinates.  If two constructed fits have
comparable scales and a common interior region containing a ball of
radius comparable to $\ell$, then, on a fixed smaller ball in that
region,
\begin{equation}
 \label{ce:overlap}
 \|D^i(f-f')\|_{L^\infty}
       \le C m_0^2\ell^{2k_0-1-i},\qquad 0\le i\le J.
\end{equation}
The constants in this last assertion can depend on the fixed scale ratio
and the fixed interior room.  No measurability of $x\mapsto f_{x,j}$ is
required.
\end{lemma}

\begin{proof}
We first give a local correction argument, keeping all its domains a
fixed positive distance apart.  Work in units of one fitting scale and
normal coordinates over a smooth minimal input $S'$.  Suppose that its
slope and a sufficiently large finite number of its scaled derivatives
are small, the mass is bounded, the selected integer is $Q$, and
$|z|\le Ca$.  Here and throughout this proof constants may incorporate
fixed factors such as $2^{k_0}$.

For the ambient normal displacement field $Z(X)=X-\pi(X)$, direct
differentiation gives
\begin{equation}
 \label{ce:normaltrace}
 \tr(P\,DZ)
  =e-P_{hh}:A_z(I-A_z)^{-1}
  =e-H_{S'}\cdot z+O(|A|^2|z|^2+|A|e|z|).
\end{equation}
Testing stationarity with $\chi^2Z$, as in the classical height-to-tilt
estimate \cite[Lemma~31]{Allard1987}, absorbing the last term, and using
$|P_{vh}|^2\le e$ gives
\[
 \int\chi^2e\,\dd\mu
 \le C\int\bigl(|\nabla\chi|^2+|A|^2\chi^2\bigr)|z|^2\,\dd\mu
 \le Ca^2.
\]
Here $H_{S'}=0$ is essential: it removes the term linear in height.
\Cref{lt:scalar} then gives $|M|\le Ca^2$ on smaller balls.  The integer
is the previously selected $Q$: moving a normalized smooth mass test
between the old and new projections changes its average by $O(a)$,
while the old mass defect is $O(a^2)$. For small $a$ its average
therefore remains within $1/4$ of $Q$. The exceptional projection counts occupy only
$O(a^2)$ of each such ball.  Apply \Cref{lt:scalar} successively through
overlapping balls if the smaller region is not concentric.

Let $u=Q^{-1}\pi_\#(z\mu)$, interpreted as a distributional normal
section using $\vol_{S'}$.  Normal stationarity and horizontal
stationarity with a linear height multiplier imply
\begin{equation}
 \label{ce:jacobi}
 \left|\int_{S'}\left(
       \nabla^\perp u:\nabla^\perp v-\tr(A_uA_v)\right)
          \dd\vol_{S'}\right|
       \le Ca^2\|v\|_{C^N}
\end{equation}
for compactly supported smooth normal fields $v$ in a smaller region.
Derivatives of $u$ in this formula are distributional.  For completeness,
the trace of the normal test $v\circ\pi$ is
\[
 -P_{hh}:A_v(I-A_z)^{-1}
       +P_{vh}:\nabla^\perp v\,(I-A_z)^{-1}.
\]
Its terms of first order in $(z,P_{vh})$ are
$-\tr(A_zA_v)+P_{vh}:\nabla^\perp v$; the constant trace is
$-H_{S'}\cdot v=0$.  The horizontal lift multiplied by $z\cdot w$
identifies $\nabla^\perp u$ with $Q^{-1}\pi_\#(P_{vh}\mu)$ to first
order.  All remainders contain $e$, $|z|\sqrt e$, or $|z|^2$, against
bounded derivatives of the test.  Their mass is $O(a^2)$.  This proves
\eqref{ce:jacobi}, including the normal-connection terms.  Its bilinear
form can also be read directly from the quadratic part of the graph
area integrand, whose differential blocks are
$(I-A_v,\nabla^\perp v)$.

Here is the local elliptic argument that turns \eqref{ce:jacobi} into an
exact minimal correction.  Choose an integer $q>m/2+1$ sufficiently
large that the right side of \eqref{ce:jacobi} is bounded by the
$H^{q+2}$ norm of $v$, and choose a positive integer $s>m/2+J+3$.
Fix a nested chain between the available radius $128$ and the required
radius $80$, for example with intermediate radii $120,112,104,96,88$,
subdividing these intervals into finitely many further radii for the
Sobolev bootstrap. Successive restrictions below use this fixed chain;
all cutoffs and extensions have a positive buffer.
For any nested pair of these working balls the Jacobi operator
$\mathcal J$ has a local right inverse $\mathcal R$ with
\begin{equation}
 \label{ce:inverse}
 \|\mathcal RF\|_{H^{t+2}}
       \le C_t\|F\|_{H^t},\qquad -q-2\le t\le s,
 \qquad \mathcal J\mathcal RF=F
 \quad\hbox{on the smaller ball}.
\end{equation}
Only finitely many coefficient derivatives are involved.  One explicit
construction removes boundary compatibility issues as follows.  Embed
the working balls in a fixed flat torus, extend the small coefficient
perturbations of $-\Delta$ by a fixed cutoff, and choose a smooth scalar
bump $b$, supported outside the larger ball, with integral one.  The
operator
$-\Delta+b\int_{\mathbb T^m}(\cdot)$ is invertible: integrate the
equation to obtain the mean of the unknown, solve the remaining
mean-zero equation by Fourier series, and restore that mean.  The
Fourier multiplier gains two derivatives on every integer Sobolev
scale.  Its sufficiently small coefficient perturbation is invertible
on each of the finitely many scales in \eqref{ce:inverse} by a Neumann
series.  Multiply $F$ by a cutoff equal to one on the desired smaller
ball and apply this inverse.  On that ball the added bump term is zero,
so the asserted equation is exact.

Apply this construction to the residual $\mathcal Ju$.  On a smaller
ball,
$v_0=u-\mathcal R(\mathcal Ju)$ satisfies $\mathcal Jv_0=0$ and
$u-v_0=O(a^2)$ in $H^{-q}$.  Moreover $u=O(a)$ in that norm, since it
is a measure of mass $O(a)$.  Interior regularity gives
$\|v_0\|_{H^{s+2}}\le Ca$ after another restriction.  This last
estimate follows directly from \eqref{ce:inverse}: apply the equation
to nested cutoffs of $v_0$; each commutator has order at most one, so
each step improves the available Sobolev order by one.  Iteration
reaches $s+2$ and then Sobolev embedding gives every fixed derivative
bound required here.

Write the minimal graph equation for a normal section as
$\mathcal Jv+\mathcal N(v)=0$.  In $H^{s+2}$ on these fixed balls,
the product rule and $s>m/2$ give
\[
 \|\mathcal N(v)\|_{H^s}\le C\|v\|_{H^{s+2}}^2,
 \quad
 \|\mathcal N(v)-\mathcal N(w)\|_{H^s}
 \le C(\|v\|_{H^{s+2}}+\|w\|_{H^{s+2}})
           \|v-w\|_{H^{s+2}}.
\]
Extend $v_0$ by a cutoff and solve
$w=-\mathcal R\mathcal N(v_0+w)$ on a ball of radius $Ca^2$ in
$H^{s+2}$, putting the nonlinear cutoff equal to one on the ultimate
working region.  This is a contraction for small $a$.  On that region
$v=v_0+w$ defines an exact minimal graph and
$v-u=O(a^2)$ in $H^{-q}$.  Interior differentiation of the equation
gives smoothness and any further fixed derivative estimates.

Converting this normal graph to a graph over $P_0$ gives the desired
centered fit.  Vertical displacement at the fixed horizontal base is
a smooth linear transformation of $z$, up to $O(|z|^2)$.  Moving the
argument of a smooth test from one projection to the other costs
$O(a^2)$ when multiplied by the height.  Finally, evaluating the graph
correction against $\pi_\#\mu$ instead of $Q\vol_{S'}$ costs
$O(a^3)$, by the mass-defect bound and its $O(a)$ size.  The first-order
term is $Q(u-v)$, proving \eqref{ce:moment}.  The same argument after
changing to the new normal coordinates proves the normal version.
The new graph has height error $Ca$ and tilt $Ca^2$; another application
of \Cref{lt:scalar} gives its mass estimate.

We apply this procedure at $j=0$ using the horizontal plane as input,
with initial errors small enough that \eqref{ce:accept} holds there.
At $j>0$ use the preceding successful fit.  Its test ball has radius
$128\ell_j$, whereas the new graph is required only on
$B_{80\ell_j}(x)$.  Thus all the fixed intermediate balls just used
fit inside the available region.  Its height error in the new scaled
units is at most $2^{k_0}m_0\ell_j^{k_0-1}$.  This proves the stated
bounds at every constructed scale, including the first failure.  The
correction at scale $\ell$ has physical derivatives of order $i$
bounded by $Cm_0\ell^{k_0-i}$.  Summing over the earlier dyadic scales
is finite for $i\le J<k_0$.  This maintains the small physical
$C^J$ bounds.  The order choices here are noncircular: first fix $J$
and the finite test order in \eqref{ce:jacobi}, then $q,s$ and the
finite coefficient order required by \eqref{ce:inverse}, then $k_0$
and finally the smallness threshold for $m_0$.  The required
coefficient order may exceed $J$.  Its uniform \emph{scaled}
derivative bounds follow by interior bootstrapping of each exact
minimal input from its small slope and scaled $C^{2,\alpha}$ bound,
using the room between its domain and the new working region.
The plane supplies this bound at initialization.  Under the next
halving of scale the second-derivative norm contracts by $1/2$ and
its H\"older seminorm by $2^{-1-\alpha}$, while the correction adds
only $Ca_j$ and the slope increments $Ca_j$ are summable.  This
closes the inductive bound before the higher-order bootstrap is
applied.  In particular the scaled $C^{q+2}$ bounds used for changes
of projection, multiplication of Sobolev tests, and the overlap
argument below do not require physical derivatives through order
$q+2$ to be included in $J$.  Every threshold and constant is thus
uniform along the construction.

It remains to prove the quadratic compatibility.  Scale a common
interior ball to unit size and put $a=m_0\ell^{k_0-1}$.  For two fits
$f_1,f_2$, the difference $d=f_1-f_2$ solves a homogeneous uniformly
elliptic linear system: integrate the derivative of the minimal-graph
flux along the segment between their gradients.  Its coefficients
have the bounded derivatives and small ellipticity perturbation
already established.  The height bounds on the support and the
projection-count conclusion of \Cref{lt:scalar} first give
$\|d\|_{L^1}\le Ca$ on an interior ball.  Indeed, on a typical normal fiber of the first fit, choose a support
point. Both fits are within $Ca$ of that point at its horizontal
coordinate. Moving back to the horizontal coordinate of the normal
base point costs $O(a)$ by the small slopes, so the difference there
is $Ca$. The exceptional base volume
is $O(a^2)$, and the small slopes bound the difference there by a fixed
constant once it is bounded at one point.  Interior estimates for this
homogeneous system therefore give $\|d\|_{C^{q+2}}\le Ca$ on a
smaller ball.

Subtract the two fixed-base moment estimates.  This gives
$\int d\zeta\,\dd((\pi_0)_\#\mu)=O(a^2)\|\zeta\|_{H^q}$,
increasing $q$ once if necessary.  To compare this with base volume,
first replace $\pi_0X$ by $\pi_0\pi_{S_1}X$ in the test $d\zeta$.
The displacement is $O(a)$ and $d\zeta$ has derivatives $O(a)$, so
the error is $O(a^2)$.  Now \Cref{lt:scalar} changes the measure to
$Q\vol_{S_1}$ with error $O(a^3)$.  Thus, writing $J_{f_1}$ for the
area density of the graph over $P_0$,
$J_{f_1}d=O(a^2)$ in a fixed negative Sobolev norm.  Multiplication by
$J_{f_1}^{-1}$ preserves that bound.  The same homogeneous interior
estimates, now starting from this negative norm, give
$\|d\|_{C^J}\le Ca^2$.  Returning to physical coordinates gives
\eqref{ce:overlap}.
\end{proof}

\newpage
\begin{lemma}[A failed fit forces separation]
 \label{ce:failed}
Fix $0<c_0<1/10$ and a tilt threshold $\varepsilon_0>0$.  By reducing
the admissible upper bound for $m_0$ in \Cref{ce:fit}, there is a
constant $c>0$ with the following property.  If the first failure at
$x$ has scale $\ell$, then in every base ball of radius $c_0\ell$
whose center lies in $B_{30\ell}(x)$, normal projection onto the failed
fit has a subset of base volume at least $c\ell^m$ on which the total
integer count is $Q$ and two entries have separation at least
$c m_0\ell^{k_0}$.  Those two entries have tilt at most
$\varepsilon_0$ relative to the fit and each has multiplicity strictly
less than $Q$.  The word ``entry'' refers to a point where the
projection Jacobian is nonzero, with its integer multiplicity; no
smoothness of a sheet is asserted.  In particular no failure occurs
when $Q=1$.
\end{lemma}

\begin{proof}
Suppose the assertion fails along a sequence $m_0\downarrow0$.  Dilate
each failed fitting ball by $\ell^{-1}$ and put
$a=m_0\ell^{k_0-1}$.  Relative to its failed fit, $|z|\le Ca$ and
$D+|M|\le Ca^2$ on fixed interior balls.  The reference graphs converge
smoothly to a plane.  Identify their bases by their graph coordinates,
including their volume densities in the transported measures.  After
taking a subsequence on the connected interior disk $B_{70}$,
\begin{equation}
 \label{ce:limits}
 a^{-2}\pi_\#(|z|^2\mu)\rightharpoonup\nu,
 \qquad a^{-2}\pi_\#(B\mu)\rightharpoonup\beta,
 \qquad a^{-2}D\rightharpoonup d,
 \qquad a^{-2}M\rightharpoonup\mathfrak m,
 \qquad d=\tr\beta.
\end{equation}
Here $\nu,d$ are positive, $\beta$ is a positive semidefinite
matrix-valued measure, and $\mathfrak m$ is signed.  The bound
$|z|/a\le C$ and $\pi_\#\mu=Q\vol_S+O(a^2)$ in total variation
imply $\nu\le C\,\dd y$ on every compact interior set.

Failure means that the fixed-base vertical error exceeds $a$ at a
support point over $B_{64}$.  Since the failed graph has uniformly small
slope, its nearest normal distance is comparable to that vertical error.
Thus $|z|\ge ca$ at this point.  The change of projection moves its base
point by $O(a)$, so a fixed interior ball about it still lies over $B_{70}$.
By \eqref{ce:normaltrace} and the vanishing curvature of the
scaled fits, the tangential Laplacian of $|z|^2$ is bounded below by
$-o(a^2)$.  More explicitly, absorb the $e|A||z|$ term in the
nonnegative $e$ term and retain $-C|A|^2|z|^2$.  Adding
$o(a^2)|X-X_1|^2$ gives a nonnegative subharmonic function on a fixed
interior ambient ball about that point.  Its mean-value inequality
from \Cref{lt:monotonicity} shows that $\nu(B_{70})\ge c>0$.

The signed estimate gives the crucial additional inequality
\begin{equation}
 \label{ce:limitsign}
 2\mathfrak m\le d.
\end{equation}
We check its physical-scale hypothesis even if $\ell$ stays bounded
away from zero.  Before applying \Cref{sx:signed}, dilate the original
physical space by $\lambda_0=m_0^\sigma$, where
$0<\sigma<1/(J-1)$ and $J$ is increased beforehand if further fixed
coordinate derivatives are required.  A physical derivative of order
$i\le J$ of a fit is now bounded by
$Cm_0\lambda_0^{1-i}$, hence uniformly bounded.  A comparable test
ball has physical radius $l\asymp\lambda_0\ell$, and its physical
height is at most $C\lambda_0\ell a$.  With $d_0=Ca^2$ and
$T=-\log d_0$, choose
\begin{equation}
 \label{ce:physicalparameters}
 0<\kappa<\min\left\{\frac\sigma2,
                          \frac1{2(k_0-1)}\right\}.
\end{equation}
Then $l+\sup|z_{\rm phys}|\le d_0^\kappa$ for all sufficiently
small $m_0$, uniformly in $0<\ell\le1$.  The fits are minimal, so
their $\Lambda$ consists only of physical height squared and
$\Lambda/a^2\le C\lambda_0^2\ell^2\pi_\#\mu\to0$ locally.
Apply \Cref{sx:signed} on finitely many interior balls to an arbitrary
nonnegative smooth test, rescaling its size and Lipschitz constant by
fixed factors.  Divide by $a^2$ and pass to the limit.  Its error is
$O(e^{-F_Q(T)})$, which tends to zero, proving \eqref{ce:limitsign}.

We next prove that $\nu$ is nonzero in \emph{every} open interior
subball.  This is a finite-radius unique-continuation argument for
the limiting measures, and requires no regularity theorem for a
multiple-valued limiting function.  Fix an interior center $y_0$ and
an integer $p\ge4$, and set
\[
 \phi_s(y)=s^{-m}\left(1-\frac{|y-y_0|^2}{s^2}\right)_+^p,
 \qquad
 \psi_s(y)=2p\,s^{-m}\left(1-\frac{|y-y_0|^2}{s^2}\right)_+^{p-1}.
\]
In this paragraph only, define
\[
 \mathcal H=\int\psi_s\,\dd\nu,
 \quad \mathcal A=s^2\int\phi_s\,\dd\mathfrak m,
 \quad \mathcal D=s^2\int\phi_s\,\dd d,
 \quad \mathcal R=\int\psi_s\,\dd\beta[y-y_0,y-y_0].
\]
Dots denote differentiation in $\log s$.  Stationarity gives
\begin{equation}
 \label{ce:limitidentities}
 \dot{\mathcal A}=2\mathcal A-\mathcal D+\mathcal R,
 \qquad \dot{\mathcal H}=2\mathcal D,
 \qquad \mathcal R\mathcal H\ge\mathcal D^2,
 \qquad 2\mathcal A\le\mathcal D.
\end{equation}
Here are the identities before the limit.  The horizontal test
$(y-y_0)\phi_s$ gives the first equation.  The normal test $z\phi_s$
identifies $s^2\int\phi_s e/a^2\,\dd\mu$ with
\[
 \int\psi_s\,\frac{z\cdot P_{vh}(y-y_0)}{a^2}\,\dd\mu+o(1).
\]
The horizontal test $(y-y_0)\psi_s|z|^2/a^2$ gives twice this flux
for the derivative of the height integral.  Terms containing both
tilt and $|z|^2/a^2$ have mass $O(a^2)$, and the geometric errors tend
to zero by \eqref{ce:normaltrace} and \Cref{lt:curvedvariation}.
Finally $P_{hv}P_{vh}=B-B^2\le B$ and Cauchy--Schwarz bound the square
of the flux by the product defining $\mathcal R\mathcal H$.
The last inequality is \eqref{ce:limitsign}.

Suppose there is a positive first radius $s_0$ at which this height
becomes nonzero.  Put $t=\log s$, $t_0=\log s_0$, and choose a fixed
$t_1>t_0$ with all the balls strictly interior.  On $(t_0,t_1]$ put
$P=\mathcal A/\mathcal H$ and $n=\mathcal D/\mathcal H$.  Then
\begin{equation}
 \label{ce:limitfrequency}
 n\ge0,\qquad 2P\le n,\qquad
 \dot P\ge 2P-n+n^2-2nP.
\end{equation}
For $P\ge1$ the last expression is nonnegative on $n\ge2P$.
Consequently $P(t)\le\max(1,P(t_1))$ going inward.
On the zero-height core, the second identity in
\eqref{ce:limitidentities}, applied also to smaller balls and other
centers, gives $d=0$.  To pass the total-variation estimate of
\Cref{lt:scalar} to this core, reuse the auxiliary physical dilation
$\lambda_0=m_0^\sigma$ above: it leaves normalized $M,D$ unchanged
and makes $a^{-2}\Lambda\to0$.  For a ball $B$ with $3B$ compactly
inside the core, take $\chi\in C_c(3B)$ equal to one on $2B$.
For every $\zeta\in C_c(B)$ with $|\zeta|\le1$, that estimate gives
\[
 |\mathfrak m(\zeta)|
 \le C\lim_i\int\chi\,a_i^{-2}(D_i+\Lambda_i)
 =C\int\chi\,\dd d=0.
\]
Hence $\mathfrak m=0$ in the core, without a boundary-continuity
assumption on the limiting measures.  The compact kernels converge
uniformly as $t\downarrow t_0$ and vanish on the boundary of the
core, so $\mathcal A(t)\to0$ even if $\mathfrak m$ charges that
boundary.  Moreover
\[
 (e^{-2t}\mathcal A)^{\displaystyle\cdot}
 =e^{-2t}(\mathcal R-\mathcal D)
 \ge-\tfrac14e^{-2t}\mathcal H,
\]
because $n^2-n\ge-1/4$.  The function $\mathcal H$ is nondecreasing.
Integrating from $t_0$ therefore gives the explicit lower bound
\[
 P(t)\ge-\frac{e^{2(t-t_0)}-1}{8}.
\]
With $P$ bounded on this finite interval,
\eqref{ce:limitfrequency} implies $\dot P\ge n^2/2-C$.
Integration gives a uniform bound for
$\int_{t_0+\epsilon}^{t_1}n^2\,\dd t$.  Cauchy--Schwarz and
$(\log\mathcal H)^{\displaystyle\cdot}=2n$ now preclude
$\mathcal H(t)\to0$ at $t_0$, a contradiction.  If $\nu$ vanished
in an open ball but were nonzero elsewhere in the connected interior
disk, choose a point of the zero set sufficiently close to its boundary.
Its distance to $\spt\nu$ is positive and smaller than the available
interior radius, producing precisely the excluded $s_0$.

We also need strong, rather than merely distributional, information
about the first moment.  Divide the first moment by $a$ and write
$U_\alpha=\pi_\#((z_\alpha/a)\mu)$.  The horizontal height-multiplier
test has, in the limiting flat coordinates, the distributional form
\begin{equation}
 \label{ce:momentcompact}
 \partial_iU_\alpha=F_{i\alpha}
                          +\partial_jE_{ij\alpha},
 \qquad \|F\|\le C,\quad \|E\|=o(1).
\end{equation}
In the flat case one may take
$F_{i\alpha}=\pi_\#(P_{v_\alpha h_i}\mu/a)$ and
$E_{ij\alpha}=\pi_\#((z_\alpha/a)B_{ij}\mu)$.
Cauchy--Schwarz bounds the first mass by $C$ and the second by
$C\int e\,\dd\mu=O(a^2)$.  In the actual coordinates, multiply
the identity by the smooth inverse base metric and volume density.
The additional terms have the same form: coefficients converge
smoothly, flux masses remain bounded, and the terms containing height
times tilt have mass tending to zero.

Take the divergence in \eqref{ce:momentcompact} and solve on a smaller
ball by a cutoff Newtonian potential.  First derivatives of the
Newtonian kernel convolved with bounded measures form a precompact
family in local $L^1$: their local kernels are integrable and have a
uniform $L^1$ translation modulus.  The second-derivative potential
of $E$ tends to zero in measure by the weak $(1,1)$ estimate proved
in \Cref{lt:potential}.  The remaining harmonic
distributions are compact on interior balls by their bounded negative
norms.  Thus $U$ is precompact locally in measure after replacing it
by its projection-count density.  This replacement is legitimate:
the area formula and the Jacobian estimate give
\[
 \left\|U-\left(\sum_{X\in\pi^{-1}(y)}
                  \theta(X)z(X)/a\right)\dd\vol_S(y)\right\|
       =o(1).
\]
One can make the potential argument for functions throughout by
mollifying with a radius chosen separately for each member of the
sequence so that the count density changes by $o(1)$ in $L^1$.
The above small total-variation difference remains $o(1)$ after
mollification.  The densities are uniformly integrable, since
$|z|/a\le C$ and their total counts differ from $Q$ by $o(1)$ in
$L^1$.  Finally \eqref{ce:moment} says that their distributional
limit is zero.  Precompactness in measure and uniform integrability
therefore show that the normalized count first moment tends to zero
in measure (indeed in local $L^1$).

On each fixed subball, the limiting second moment is positive.
Fibers of total count different from $Q$ have vanishing count mass by
\Cref{lt:scalar}; points with tilt exceeding a number tending slowly
to zero have vanishing mass by \eqref{ce:fitenergy}.  Discard both.
On a remaining fiber list its $Q$ normalized heights with repetitions
as $w_1,\ldots,w_Q$, and set $\bar w=Q^{-1}\sum_iw_i$.  The identity
\[
 \sum_{i=1}^Q|w_i|^2
   =\sum_{i=1}^Q|w_i-\bar w|^2+Q|\bar w|^2
\]
and $|w_i|\le C$ show that a positive amount of the second moment
is variance.  Hence a subset of fixed positive base volume has
variance bounded below, and therefore has two heights separated by
a fixed positive amount.  Their distinct spatial entries each have
multiplicity less than $Q$.  Restoring lengths gives the asserted
separation.  For $Q=1$ the variance is identically zero, so failure
is impossible.  If uniform constants for the prescribed comparable
subballs failed, their centers and radii would have a convergent
subsequence in the compact permitted range; the preceding argument
on an interior subball of its limit gives the same contradiction.
\end{proof}

\subsection{The center and weighted estimates}

\begin{proposition}[The glued center and its influence region]
 \label{ce:center}
Fix the desired output order $J\ge6$, apply \Cref{ce:fit} with
derivative order $J+1$, and then choose its corresponding $k_0$ and
initial parameters.  There is a $C^J$ graph $S$ on an
interior neighborhood of $0$, with $C^J$ norm at most $Cm_0$, through
$0$ and tangent to $P_0$ there, and a nonnegative Lipschitz function
$t$ on the fixed base, such that
\begin{equation}
 \label{ce:geometry}
 |H_S|+|\nabla H_S|\le Cm_0^2t^{2k_0-4},
 \qquad
 |z|\le Cm_0t^{k_0}\quad\hbox{on }\spt\mu.
\end{equation}
In the first inequality $t$ is evaluated at the horizontal coordinate
of the point of $S$; in the second it may be evaluated either at
$\pi_0X$ or at $\pi_0\pi_SX$, with a change of constant.
Writing $\mathcal I=\{t>0\}$, one has
\begin{equation}
 \label{ce:influence}
 t(0)=0,\qquad
 \sup_{B_r}t=o(r),\qquad
 |\mathcal I\cap B_r|=o(r^m).
\end{equation}
In particular the support height is zero off the influence region.
On $S$ itself its normal-projection version has the same vanishing
volume density.  If $\rho$ is distance to $0$ on $S$, then
$|z|\le C\rho^2$ on the support in a sufficiently small tube.
\end{proposition}

\begin{proof}
Let $\ell_x$ be the first failed scale at $x\in B_5$, and put
$\ell_x=0$ if there is no failure.  Define
\begin{equation}
 \label{ce:tent}
 t(y)=\sup_{x\in B_5}(\ell_x-|x-y|)_+,
 \qquad y\in B_2.
\end{equation}
Every function in this supremum is $1$-Lipschitz, so $t$ is
$1$-Lipschitz and $\mathcal I$ is the union of the failure balls.

We describe the gluing before proving the density properties.  For
$\epsilon>0$ use $t_\epsilon=\max(t,\epsilon)$.  Choose a covering
by balls of radius $c t_\epsilon(y_i)$, with $c>0$ a sufficiently
small fixed constant, by first taking disjoint balls of a smaller
such radius and then enlarging them.  The Lipschitz bound makes
radii of intersecting enlarged balls comparable; the covering has
bounded overlap.  Smooth bumps on these balls, divided by their
sum, give a partition $\{\chi_i\}$ with
$|D^j\chi_i|\le C_jt_\epsilon(y_i)^{-j}$.
At $y_i$ choose a dyadic constructed scale $L_i$ between
$C_0t_\epsilon(y_i)$ and $2C_0t_\epsilon(y_i)$, capped at the initial
scale $1$.  The constant $C_0$ is fixed and the covering radii can be
reduced to leave all required overlap room.  Since
$\ell_{y_i}\le t(y_i)$, these scales have been constructed.
Where the cap is active, $L_i\asymp t_\epsilon(y_i)$ still holds,
with a constant depending on $C_0$.  Let
$f_\epsilon=\sum_i\chi_i f_{y_i,L_i}$.

On an overlapping cell of size $L$, fix one of its fits $f_*$ and
write $f_\epsilon-f_*=\sum_i\chi_i(f_i-f_*)$.  By
\eqref{ce:overlap}, every derivative of order $j\le J+1$ in this
identity is at most $Cm_0^2L^{2k_0-1-j}$, by the choice of derivative
order $J+1$ in \Cref{ce:fit}.  Each $f_i$ is minimal.
The mean-curvature operator and its first derivative involve at
most three derivatives of the graph.  Substitution therefore gives
\[
 |H_{\operatorname{graph}f_\epsilon}|+
 |\nabla H_{\operatorname{graph}f_\epsilon}|
       \le Cm_0^2t_\epsilon^{2k_0-4}.
\]
The support height is at most $Cm_0t_\epsilon^{k_0}$, by the height
bounds for the fits.  Uniform $C^{J+1}$ bounds give, along a sequence
$\epsilon\downarrow0$, a local $C^J$ limit $f$.  Passing to the
limit proves \eqref{ce:geometry}, including at $t=0$.  Uniform small
slope supplies a normal tube on each smaller chart.  The change in
horizontal argument under nearest projection is at most
$Cm_0|z|\le Cm_0^2t^{k_0}$, which is a small fraction of $t$ when
$t>0$.  The Lipschitz bound on $t$ therefore proves the assertion
about the projected argument, including the zero set by a limit.

We next use multiplicity defects to control the tents.  Choose the
comparable subball in \Cref{ce:failed} strictly inside the original
failure ball.  Small slope and the much smaller height displacement
then imply, in the fixed horizontal projection,
\begin{equation}
 \label{ce:charge}
 \lambda(B_\ell(x))\ge c\ell^m
       \quad\hbox{for every failure }(x,\ell).
\end{equation}
This charges inside the original ball, a fact needed in the disjoint
covering below.  The change from normal to horizontal projection has
a uniformly bounded Jacobian and displaces the chosen smaller
subball by much less than its interior margin.

Fix a large finite $K$, greater than every fixed enlargement factor
that will occur in the following local arguments.  At initialization
choose $\eta<c/(K+1)^m$ in \eqref{ce:initialdefect}.  If
$|x|\le K\ell$, then
\[
 c\ell^m\le\lambda(B_\ell(x))
       \le\lambda(B_{|x|+\ell}(0))
       \le\eta(K+1)^m\ell^m,
\]
a contradiction.  The enclosing ball is inside the fixed chart;
one uses its actual radius $|x|+\ell$, not an unnecessarily larger
ball.  Thus $|x|>K\ell$ for all relevant failures, and $t(0)=0$.
If $B_{A\ell}(x)$ meets $B_r(0)$, for a fixed $A<K$, then
\[
 \ell\le\frac r{K-A},\qquad
 B_\ell(x)\subset B_{C_A r}(0),\qquad
 C_A=1+\frac{A+1}{K-A}.
\]
Combining \eqref{ce:charge} with the little-$o$ part of
\eqref{ce:initialdefect} now gives
\begin{equation}
 \label{ce:smallfailure}
 \sup\{\ell:B_{A\ell}(x)\cap B_r(0)\ne\varnothing,
                      \ (x,\ell)\hbox{ a failure}\}=o(r).
\end{equation}
All limits here hold for the one fixed choice of the initial data.
This first excludes large intrusions using a uniform bound and only
then uses differentiation to obtain the little-$o$ conclusion.

For small $r$, all original failure balls meeting $B_r$ lie in
$B_{2r}$ by \eqref{ce:smallfailure}.  Select disjoint original balls
$B_i$ whose fivefold enlargements cover their union.  The charges in
\eqref{ce:charge} are disjoint, so
\[
 |\mathcal I\cap B_r|
       \le C\sum_i|B_i|
       \le C\sum_i\lambda(B_i)
       \le C\lambda(B_{2r})=o(r^m).
\]
The same small-failure estimate gives $\sup_{B_r}t=o(r)$.
The support contains $0$, so its zero height at $t(0)=0$ shows that
$S$ passes through $0$.  The superlinear height bound and
\eqref{ce:typical} force $T_0S=P_0$.  Smooth change of coordinates
transfers the volume statement to $S$, and the projected height
bound gives $|z|\le Cm_0o(\rho)^{k_0}\le C\rho^2$.
\end{proof}

We state the weighted consequences with the physical units retained.
Let $\pi=\pi_S$, $z=X-\pi(X)$, and let $\rho$ be geodesic distance
on $S$ to $0$ inside a normal coordinate ball. In these weighted
estimates and the final height argument, $B_r$ denotes the geodesic
ball $\{y\in S:\rho(y)<r\}$. Local covering balls may be taken
in fixed smooth charts; their comparison with metric balls changes
only fixed constants. Choose a fixed integer
$p\ge4$, increasing it if needed in subsequent compact-kernel tests,
and put
\begin{equation}
 \label{ce:kernel}
 \phi_r=(1-\rho^2/r^2)_+^p,
 \qquad \psi_r=2p(1-\rho^2/r^2)_+^{p-1},
 \qquad \mathbf u=\nabla_S\rho,\quad K_z=(I-A_z)^{-1}.
\end{equation}
Weighted integrals from this point onward include $r^{-m}$; when
there is any possibility of confusion we write this factor explicitly.
The heights $z$ are not divided by $r$.  The symbols $B,e,M,D,L,\Lambda$
now always refer to this center.  Define
\begin{equation}
 \label{ce:radialheight}
 \begin{aligned}
 b&=B[K_z\mathbf u,K_z\mathbf u],
 &h_r&=1-\frac b{|K_z\mathbf u|^2},\\
 H(r)&=r^{-m}\int\psi_r|z|^2h_r\,\dd\mu,
 &\delta(r)&=r^{-1}\sup\{|z|:X\in\spt\mu,\ \rho(\pi X)<r\}.
 \end{aligned}
\end{equation}
In particular $H$ already contains the radial projection factor.
One has $0\le h_r\le1$, since $0\le B\le I$.  At a fixed support
point this factor has no dependence on the radius $r$; the subscript
only distinguishes it from other height variables.  On the central
fiber it may be assigned any value, because $z=0$ there.

\begin{proposition}[Weights up to the boundary of the ball]
 \label{ce:weights}
Either the support lies on $S$ in a neighborhood of $0$, or there
exist $r_0>0$, $c,C>0$, and a finite $J_0$ such that, for every
$0<r<r_0$,
\begin{equation}
 \label{ce:height}
 0<\delta(r)\le Cr,
 \qquad c\delta(r)^{J_0}\le\frac{H(r)}{r^2}
                         \le C\delta(r)^2.
\end{equation}
In the latter alternative, throughout $\rho<r$ one also has
\begin{equation}
 \label{ce:pointwisecurvature}
 |H_S|+|\nabla H_S|\le C\delta(r)r,
\end{equation}
and
\begin{equation}
 \label{ce:curvature}
 r^{-m}\int\phi_r\bigl(|H_S|^2+|\nabla H_S|^2\bigr)\,\dd\mu
       \le C r^{-m}\int\phi_r|z|^2\,\dd\mu.
\end{equation}
The constants may depend on the fixed center construction.  One may
take any fixed $J_0\ge2+(m+p-1)/k_0$ after changing $c$.
\end{proposition}

\begin{proof}
We give the geometric estimate that pays for points close to the
boundary $\rho=r$.  If $t(y)>0$, choose a failure $(x,\ell)$ for
which
\begin{equation}
 \label{ce:witness}
 \ell-|x-y|>t(y)/2.
\end{equation}
Then $\ell>t(y)/2$, $|x-y|<\ell$, and the Lipschitz bound gives
$t(v)\le(2+A)\ell$ when $|v-y|\le A\ell$.  An upper comparison
of $\ell$ with $t(y)$ is neither asserted nor needed.  On this whole
neighborhood the support has height at most $Cm_0\ell^{k_0}$ over
$S$, and its mass over any ball of radius comparable to $\ell$ is
at most $C\ell^m$ by the ambient upper mass bound.

In each fixed small comparable subball within the permitted distance
from $x$, \Cref{ce:failed} gives mass at least $c\ell^m$ with
distance from $S$ at least $cm_0\ell^{k_0}$ and with tilt relative
to $S$ less than, say, $1/4$.  To verify this transfer, the normal
fiber of the failed fit is transverse to the small-slope graph $S$.
It has only one intersection with that graph in the working tube,
and distance to $S$ is comparable to distance along the fiber to
that intersection.  Of two separated entries at least one is a
fixed fraction of their separation from $S$.  Their tilt relative
to $S$ is still small, since both reference graphs have small slope.
Restrict the subball slightly first: all changes of projection
argument are $o(\ell)$ and so preserve its interior margin.  On
this mass $h_r\ge3/4$ after reducing the tilt threshold in
\Cref{ce:failed} once.

For $y$ with $\rho(y)<r$, the witnessing scale satisfies $\ell=o(r)$
by \eqref{ce:smallfailure}.  There is a ball of radius $c_1\ell$
at distance $O(\ell)$ from $y$, lying inside $\rho<r$, on which
\begin{equation}
 \label{ce:payingweights}
 \inf\phi_r\ge c\sup_C\phi_r,
 \qquad \inf\psi_r\ge c\sup_C\psi_r,
 \qquad \inf\psi_r\ge c(\ell/r)^{p-1}.
\end{equation}
Here $C$ is a cell of side $c_2\ell$ containing $y$, and the
suprema concern its intersection with $\rho<r$.  If $y$ is already
a fixed multiple of $\ell$ inside the boundary, take a nearby
subball there.  Otherwise move its center inward along the radial
geodesic by $C_1\ell$ and take $c_1$ much smaller than $C_1$.
The boundary distance on the new ball is at least a fixed multiple
of the largest boundary distance on the cell and at least
$c\ell$.  Since
$1-\rho^2/r^2$ is comparable to $(r-\rho)/r$ near the boundary,
this gives \eqref{ce:payingweights}; in the central half ball all
weights are bounded below.  The displacement is bounded by a
fixed small multiple of $\ell$, so the new ball is among those
allowed in \Cref{ce:failed}.  These assertions are unchanged by
the smooth base-coordinate conversions.  Thus they also apply to
cells straddling $\rho=r$.

\begin{figure}[tb]
\centering
\begin{tikzpicture}[x=1.55cm,y=1.55cm,>=Latex,font=\small]
 \begin{scope}
  \clip (-.25,-1.1) rectangle (4.4,1.65);
  \fill[blue!5] (0,0) circle (3);
 \end{scope}
 \fill (0,0) circle (1.4pt) node[left] {$0$};
 \draw[->,gray] (0,0)--(3,0) node[midway,below] {$r$};
 \node[blue!45!black] at (1.15,1.34) {base ball $B_r$};
 \draw[red!65!black,dashed,thick] (3.08,.65) circle (.85);
 \fill[red!65!black] (3.08,.65) circle (1.4pt);
 \draw[red!65!black] (3.08,.65)--(3.89,.91)
      node[midway,above] {$\ell$};
 \node[red!65!black,anchor=west] at (3.57,1.51) {failure ball};
 \draw[red!65!black,->] (3.6,1.39)--(3.57,1.25);
 \draw[gray,fill=gray!12] (2.63,.3) rectangle (3.18,.9);
 \begin{scope}
  \clip (-.25,-1.1) rectangle (4.4,1.65);
  \draw[blue!65!black,thick] (0,0) circle (3);
 \end{scope}
 \node[blue!65!black,anchor=west] at (3.02,-.33) {$\rho=r$};
 \fill (2.78,.58) circle (1.4pt) node[above left] {$y$};
 \fill[green!14] (1.73,.58) circle (.32);
 \draw[green!40!black,thick] (1.73,.58) circle (.32);
 \draw[->,thick,orange!75!black] (2.58,.58)--(2.12,.58);
 \node[green!35!black,align=center,anchor=north] at (1.45,-.32)
      {inward subball\\with positive height mass};
 \draw[green!35!black,->] (1.64,-.29)--(1.73,.2);
 \node[anchor=north,align=center] at (3.81,-.71)
      {cell crossing\\the kernel boundary};
 \draw[gray,->] (3.66,-.65)--(3.14,.3);
\end{tikzpicture}
\caption{An inward subball controls the boundary weights on the reference center. A failure at
scale $\ell$ supplies separated projection entries on every sufficiently small comparable
subball in its comparison region. Moving such a subball inward by $O(\ell)$
gives positive height mass where the polynomial kernel has a controlled lower
bound. This pays for a local cell even when that cell crosses $\rho=r$. The
picture is schematic; no circular shape of a failure region is assumed.}
\label{fig:collar}
\end{figure}


The separated mass on this inward ball gives
\begin{equation}
 \label{ce:payingheight}
 \delta(r)r\ge c m_0\ell^{k_0}.
\end{equation}
If the support fails to lie on $S$ in every neighborhood, then
$\delta(r)>0$ for every small $r$.  Choose a point with height at
least half the defining supremum.  Its height is nonzero, so it
has a witness as in \eqref{ce:witness}; \eqref{ce:geometry} and
\eqref{ce:payingheight} give
\[
 c m_0\ell^{k_0}\le\delta(r)r
                         \le C m_0\ell^{k_0}.
\]
The inward ball has mass at least $c\ell^m$, height at least
$cm_0\ell^{k_0}$, and $h_r\ge3/4$.  Consequently
\[
 \frac{H(r)}{r^2}
 \ge c\delta(r)^2(\ell/r)^{m+p-1}
 \ge c(m_0)\delta(r)^{2+(m+p-1)/k_0}.
\]
In the last step use
$\ell/r\ge c(m_0)\delta(r)^{1/k_0}r^{1/k_0-1}$ and $r\le1$.
The upper bound in \eqref{ce:height} follows from bounded normalized
mass and $h_r\le1$.  The bound $\delta(r)\le Cr$ follows from
\Cref{ce:center}.  This proves \eqref{ce:height}.

For any point with $t>0$, its witness and
\eqref{ce:payingheight} give
\[
 |H_S|+|\nabla H_S|
 \le Cm_0^2\ell^{2k_0-4}
 \le Cm_0\ell^{k_0-4}\delta(r)r
 \le C\delta(r)r.
\]
At $t=0$ the left side is zero.  This proves
\eqref{ce:pointwisecurvature}.

Finally group points with $t>0$ according to the dyadic scale of
a witnessing failure.  At scale $\ell=2^{-j}$ cover the group by
grid cells of side $c_2\ell$, taking each grid cell only once.
For each occupied cell choose a witness from that group and an
inward ball as above.  On the cell the curvature-square integrand
is at most $Cm_0^4\ell^{4k_0-8}$.  Its mass is at most
$C\ell^m$.  On the corresponding inward ball the integral of
$\phi_r|z|^2$ is at least
$c m_0^2\ell^{2k_0+m}\sup_C\phi_r$ by
\eqref{ce:payingweights}.  Hence
\[
 \int_C\phi_r\bigl(|H_S|^2+|\nabla H_S|^2\bigr)\,\dd\mu
 \le Cm_0^2\ell^{2k_0-8}
                  \int_{B_C}\phi_r|z|^2\,\dd\mu.
\]
The displacement from a cell to its inward ball is $O(\ell)$,
so these balls have bounded overlap at each fixed dyadic scale.
Sum first over the cells and then over $j$.  The geometric series
$\sum_j2^{-j(2k_0-8)}$ converges because $k_0>4$.  Points with
$t=0$ contribute nothing, and division by $r^m$ gives
\eqref{ce:curvature}.
\end{proof}

\begin{corollary}[Regularity when height vanishes]\label{ce:constancy}
If the support lies on the center $S$ in a neighborhood of the origin,
then the origin is regular.
\end{corollary}

\begin{proof}
Write the restricted varifold as $\vartheta\,\HH^m\llcorner S$,
allowing the integer density $\vartheta$ to be zero off its original
support. This is the standard constancy argument; compare
\cite[Chapter~8, Theorem~4.1]{Simon2018}. Tangential first variation gives
$\int_S\vartheta\,\Div_S X=0$ for all compactly supported tangent
fields.  Thus $\vartheta$ is constant almost everywhere on a connected
smaller disk.  The constant is $Q$, by \eqref{ce:typical}.  Normal
first variation gives $H_S=0$, and the small-gradient minimal graph
equation bootstraps $S$ to a smooth graph.  The origin is therefore
regular.
\end{proof}

It remains to exclude the positive-height alternative, which is the
purpose of the inward comparison.
